% FILE: chapters/ch01_s2_dev_environment.tex
\section{Development Environment on macOS}\label{sec:dev_environment_macos}

Before building machine learning algorithms or deploying edge firmware to microcontrollers, you need a stable, professional development environment. Many beginners fail in machine learning not because the mathematics is too difficult, but because their computer environment breaks: packages conflict, paths are misconfigured, Python versions collide, and code that ran yesterday suddenly throws cryptic import errors today.

On macOS, this problem is especially common because Apple ships macOS with a built-in ``system Python'' that the operating system itself relies on for internal system tasks. Modifying or polluting this system Python can destabilize your operating system.

In this section, we build a clean, resilient, and professional machine learning workspace on macOS from the ground up:
\begin{enumerate}
    \item \textbf{The Terminal}: Conquering the command line with zero prior experience.
    \item \textbf{Package Management with Homebrew}: Installing an isolated, modern Python runtime safely.
    \item \textbf{Isolated Environments}: Mastering Python's built-in \code{venv} and Miniconda to prevent dependency conflicts.
    \item \textbf{Package Management with \code{pip}}: Installing libraries and freezing dependencies with \code{requirements.txt}.
    \item \textbf{Interactive Exploration}: Running exploratory data workflows in Jupyter Notebook and JupyterLab.
    \item \textbf{Professional Editor}: Configuring Visual Studio Code (VS Code) with the official Python extension.
\end{enumerate}

\begin{figure}[htbp]
\centering
\begin{tikzpicture}[
    scale=0.85,
    box/.style={rectangle, draw=keyideablue, fill=keyideablue!10, rounded corners=2mm, thick, text width=2.4cm, align=center, minimum height=1.1cm, font=\scriptsize\sffamily},
    arrow/.style={-{Stealth[length=2.5mm]}, thick, draw=keyideablue}
]
    \node[box] (brew) {1. Homebrew\\(\code{brew install})};
    \node[box, right=0.6cm of brew] (python) {2. Clean Python\\(\code{python3.11})};
    \node[box, right=0.6cm of python] (venv) {3. Isolated Venv\\(\code{.venv})};
    \node[box, right=0.6cm of venv] (pip) {4. Pinned Pip\\(\code{requirements.txt})};
    
    \node[box, below=1.0cm of pip] (jupyter) {5. JupyterLab\\(EDA \& Visuals)};
    \node[box, left=0.6cm of jupyter] (vscode) {6. VS Code IDE\\(Python Ext, Debugger)};
    \node[box, left=0.6cm of vscode] (iot) {7. TinyML Tools\\(\code{esptool}, Serial)};
    \node[box, left=0.6cm of iot] (ready) {8. Robust Workspace\\(Ready for ML)};

    \draw[arrow] (brew) -- (python);
    \draw[arrow] (python) -- (venv);
    \draw[arrow] (venv) -- (pip);
    \draw[arrow] (pip) -- (jupyter);
    \draw[arrow] (jupyter) -- (vscode);
    \draw[arrow] (vscode) -- (iot);
    \draw[arrow] (iot) -- (ready);
\end{tikzpicture}
\caption{The modular developer workspace hierarchy on macOS: building from hardware tools up to interactive IDEs.}\label{fig:dev_env_stack}
\end{figure}

% ==============================================================================
% 1. TERMINAL ESSENTIALS
% ==============================================================================
\subsection{The Terminal for Complete Beginners}

\begin{keyidea}[What is the Terminal?]
The \textbf{terminal} (also called the command line or shell) is a text-based interface to your operating system. Instead of clicking icons with a mouse, you type concise text commands to create files, launch programs, inspect hardware ports, and run machine learning scripts.
\end{keyidea}

On macOS, press \textbf{Command ($\mathbf{\cmd}$) + Space}, type \texttt{Terminal}, and press \textbf{Enter}. You will be greeted by a prompt ending with a \code{\%} symbol (the default \code{zsh} shell prompt).

\begin{analogy}[The GUI vs. The Command Line]
Using a Graphical User Interface (GUI) like macOS Finder is like dining at a restaurant with a visual menu: you can only point to what is pictured on the page. Using the command line is like walking directly into the kitchen and speaking to the chef: you have complete, precise, scriptable control over every ingredient and pan in the building.
\end{analogy}

\paragraph{The Seven Essential Navigation Commands}
Every developer must know these seven fundamental terminal commands:

\begin{table}[htbp]
\centering
\small
\begin{tabular}{lll}
\toprule
\textbf{Command} & \textbf{Full Meaning} & \textbf{Operational Action} \\
\midrule
\code{pwd} & Print Working Directory & Shows the exact folder path where you currently stand. \\
\code{ls -la} & List Files (All, Detailed) & Lists all files in the current folder, including hidden dotfiles. \\
\code{cd <folder>} & Change Directory & Moves into the named folder (e.g., \code{cd Developer/ML}). \\
\code{cd ..} & Change Directory Up & Moves up one folder level toward the root drive. \\
\code{mkdir <name>} & Make Directory & Creates a brand-new folder. \\
\code{touch <file>} & Create Empty File & Creates a new empty file (e.g., \code{touch train.py}). \\
\code{rm -i <file>} & Remove File (Interactive) & Safely deletes a file after prompting for confirmation. \\
\bottomrule
\end{tabular}
\caption{Essential macOS shell navigation commands.}\label{tab:shell_commands}
\end{table}

\paragraph{Paths: Absolute vs. Relative}
\begin{itemize}
    \item \textbf{Absolute Path}: Starts from the very root of your hard drive (\code{/}). Example: \code{/Users/alex/Developer/ML/train.py}.
    \item \textbf{Relative Path}: Starts from where your terminal currently stands. If you are in \code{/Users/alex}, the relative path to the same file is \code{Developer/ML/train.py}.
    \item \textbf{The Tilde Shortcut (\code{\~})}: Represents your user home directory (\code{/Users/<your_username>}). Typing \code{cd ~} always returns you home safely.
\end{itemize}

% ==============================================================================
% 2. HOMEBREW & CLEAN PYTHON
% ==============================================================================
\subsection{Installing Python the Right Way with Homebrew}

\begin{keyidea}[Why Homebrew?]
\textbf{Homebrew} is the missing package manager for macOS. It installs open-source software, compilers, and programming languages into isolated user directories (under \code{/opt/homebrew} on Apple Silicon Macs), completely separated from Apple's internal system files.
\end{keyidea}

\paragraph{Why You Must Never Use macOS System Python}
macOS ships with a system-level Python interpreter located at \code{/usr/bin/python3}. 
\begin{warning}[Never Touch System Python]
Never run \code{sudo pip install} into \code{/usr/bin/python3}! Apple's internal desktop automation, Xcode toolchains, and security services depend on exact, unmodified system Python libraries. Overwriting these libraries can corrupt macOS system components. Always install a separate, user-owned Python interpreter via Homebrew.
\end{warning}

\paragraph{Step-by-Step Installation}
\begin{enumerate}
    \item \textbf{Install Homebrew}: Paste the following official command into your terminal:
\begin{lstlisting}[language=bash]
/bin/bash -c "$(curl -fsSL https://raw.githubusercontent.com/Homebrew/install/HEAD/install.sh)"
\end{lstlisting}

    \item \textbf{Install Modern Python (Python 3.11 or 3.12)}:
\begin{lstlisting}[language=bash]
brew install python@3.11
\end{lstlisting}

    \item \textbf{Verify Installation}:
\begin{lstlisting}[language=bash]
which python3.11
python3.11 --version
\end{lstlisting}
\begin{lstlisting}[style=outputstyle]
/opt/homebrew/bin/python3.11
Python 3.11.8
\end{lstlisting}
The output path confirms that the interpreter resides in \code{/opt/homebrew/bin/}, completely isolated from Apple's system files.
\end{enumerate}

% ==============================================================================
% 3. VIRTUAL ENVIRONMENTS
% ==============================================================================
\subsection{Virtual Environments: \code{venv} and Miniconda}

\begin{keyidea}[The Dependency Collision Problem]
Suppose Project A (trained two years ago) requires \code{scikit-learn==0.24.0}, while your new TinyML Project B requires \code{scikit-learn==1.4.0}. If you install libraries globally, one project will inevitably crash. A \textbf{virtual environment} is a self-contained folder containing its own isolated Python interpreter and library directory.
\end{keyidea}

\begin{analogy}[The Mechanic's Dedicated Toolboxes]
A master mechanic does not throw surgical tools, heavy truck wrenches, and delicate watchmaker screwdrivers into one giant bucket. They maintain separate, dedicated toolboxes for each vehicle. If a tool in the truck toolbox breaks, the watchmaker's tools remain pristine. A virtual environment is a dedicated toolbox for a specific machine learning project.
\end{analogy}

\paragraph{Option A: Python's Built-In \code{venv} (Recommended for Most Projects)}
Python comes with a built-in virtual environment module called \code{venv}. It requires zero additional downloads.

\begin{enumerate}
    \item \textbf{Navigate to your project folder}:
\begin{lstlisting}[language=bash]
cd ~/Developer/ML
mkdir iot_vibration_project && cd iot_vibration_project
\end{lstlisting}

    \item \textbf{Create the virtual environment}:
\begin{lstlisting}[language=bash]
python3.11 -m venv .venv
\end{lstlisting}
This creates a hidden directory named \code{.venv} containing a private copy of the Python interpreter and a blank \code{site-packages} library directory.

    \item \textbf{Activate the environment}:
\begin{lstlisting}[language=bash]
source .venv/bin/activate
\end{lstlisting}
Your terminal prompt changes to display \code{(.venv)}, indicating that the isolated environment is active! Any package you install now will stay inside \code{.venv}.

    \item \textbf{Deactivate when finished}:
\begin{lstlisting}[language=bash]
deactivate
\end{lstlisting}
\end{enumerate}

\paragraph{Option B: Miniconda / Conda (Recommended for Complex C/C++ Extensions)}
While \code{venv} manages Python libraries, some machine learning packages (like PyTorch with Metal Performance Shaders or specialized TinyML compilers) require compiled C/C++ and Fortran binaries. \textbf{Conda} is a cross-platform package and environment manager that handles both Python packages and non-Python system libraries.

\begin{lstlisting}[language=bash]
# Install Miniconda via Homebrew
brew install --cask miniconda

# Create a clean, isolated environment named 'tinyml'
conda create -n tinyml python=3.11 -y

# Activate the conda environment
conda activate tinyml

# List all existing environments
conda env list
\end{lstlisting}

% ==============================================================================
% 4. PIP AND REQUIREMENTS.TXT
% ==============================================================================
\subsection{Package Management with \code{pip} and \code{requirements.txt}}

Once your virtual environment is active, you install third-party libraries using \textbf{\code{pip}} (the Python Package Installer).

\paragraph{Installing Core Data Science Packages}
\begin{lstlisting}[language=bash]
# Always upgrade pip first inside your fresh environment
pip install --upgrade pip

# Install foundational machine learning and sensor analysis packages
pip install numpy pandas scikit-learn matplotlib seaborn jupyterlab
\end{lstlisting}

\paragraph{The Gold Standard of Reproducibility: \code{requirements.txt}}
How do you ensure a teammate or an automated cloud training server can run your exact code without version mismatches? You pin your environment's exact versions into a \code{requirements.txt} file:

\begin{lstlisting}[language=bash]
# 1. Export exact pinned versions to requirements.txt
pip freeze > requirements.txt

# 2. Inspect the generated file
head -n 5 requirements.txt
\end{lstlisting}

\begin{lstlisting}[style=outputstyle]
numpy==1.26.4
pandas==2.2.1
scikit-learn==1.4.1.post1
matplotlib==3.8.3
seaborn==0.13.2
\end{lstlisting}

When sharing your project on GitHub or deploying to an edge server, the other engineer can replicate your entire software stack with a single command:
\begin{lstlisting}[language=bash]
pip install -r requirements.txt
\end{lstlisting}

% ==============================================================================
% 5. JUPYTER NOTEBOOK & JUPYTERLAB
% ==============================================================================
\subsection{Interactive Exploration: Jupyter Notebook and JupyterLab}

\begin{keyidea}[Jupyter Architecture]
A \textbf{Jupyter Notebook} (\code{.ipynb}) is an interactive computational document that combines runnable blocks of Python code (cells) with rich text, mathematical LaTeX formulas, and interactive plots.
\end{keyidea}

Unlike standard Python scripts that run from line 1 to the end and immediately terminate, a Jupyter Notebook runs an active background process called the **Kernel**. When a code cell finishes executing, all variables, arrays, and trained models remain alive in your computer's RAM, allowing you to experiment interactively without re-loading multi-gigabyte datasets each time.

\paragraph{Launching JupyterLab}
From your activated terminal environment, launch the web interface:
\begin{lstlisting}[language=bash]
source .venv/bin/activate
jupyter lab
\end{lstlisting}
Your default web browser opens to \code{http://localhost:8888/lab}, displaying the interactive notebook workspace.

\begin{warning}[The Jupyter State Pitfall: Out-of-Order Execution]
Because Jupyter lets you run cells in any arbitrary order, you can easily create an invisible, non-reproducible state. For instance, if you delete a cell that defined variable \code{x = 10}, the variable \code{x} still exists in the background kernel RAM! Before sharing any notebook, always select:
\begin{center}
\textbf{Kernel $\to$ Restart Kernel and Run All Cells}
\end{center}
If the notebook runs from top to bottom without errors, your results are genuine and reproducible.
\end{warning}

% ==============================================================================
% 6. VS CODE SETUP
% ==============================================================================
\subsection{Professional Editor: Visual Studio Code}

While Jupyter is great for data exploration, real-world machine learning pipelines, unit tests, and embedded microcontroller firmware are authored as modular Python (\code{.py}) and C/C++ (\code{.c}, \code{.h}) files. The industry standard tool is \textbf{Visual Studio Code (VS Code)}.

\paragraph{Installation and Setup}
\begin{enumerate}
    \item Install VS Code via Homebrew:
\begin{lstlisting}[language=bash]
brew install --cask visual-studio-code
\end{lstlisting}

    \item \textbf{Install the Official Python Extension}:
    Open VS Code, press \textbf{Command ($\mathbf{\cmd}$) + Shift + X} to open the Extensions tab, search for \code{Python} (authored by Microsoft), and click \textbf{Install}.

    \item \textbf{Select Your Virtual Environment Interpreter}:
    Press \textbf{Command ($\mathbf{\cmd}$) + Shift + P} to open the Command Palette, type:
    \begin{center}
    \texttt{Python: Select Interpreter}
    \end{center}
    Choose your project's local environment: \code{./.venv/bin/python}. VS Code will now provide intelligent autocompletion, real-time type checking, and integrated debugging using your isolated project packages!
\end{enumerate}

\begin{iotbox}[Edge Tooling from macOS Terminal: ESP32 \& Serial]
For IoT and TinyML development on macOS, you must communicate with microcontrollers over USB-to-UART serial bridges. Install these essential command-line tools into your environment:
\begin{lstlisting}[language=bash]
# Install ESP32 flashing and serial monitoring utilities
pip install esptool pyserial

# Detect connected ESP32 / Arduino hardware ports on macOS
ls /dev/cu.usbserial* /dev/cu.wchusbserial*
\end{lstlisting}
If connected properly, your microcontroller appears as a serial device (e.g., \code{/dev/cu.usbserial-0001}), ready for firmware flashing directly from your macOS terminal!
\end{iotbox}

\begin{takeaways}
\begin{itemize}
    \item Never install packages into macOS system Python; use Homebrew to maintain an isolated Python interpreter.
    \item Always create an active virtual environment (\code{python3.11 -m venv .venv}) for every distinct machine learning project.
    \item Pin library dependencies with \code{pip freeze > requirements.txt} to guarantee bit-for-bit reproducibility across machines.
    \item Use JupyterLab for interactive exploratory data analysis, but always test reproducibility with ``Restart Kernel and Run All''.
    \item Link VS Code directly to your project's virtual environment interpreter for seamless syntax checking and debugging.
\end{itemize}
\end{takeaways}

% ==============================================================================
% 7. CHECK YOURSELF
% ==============================================================================
\subsection{Check Yourself}

\begin{enumerate}
    \item \textbf{Environment Isolation}: Why is running \code{sudo pip install scikit-learn} on macOS considered dangerous? What can happen to your operating system?

    \item \textbf{Virtual Environment Mechanics}: You create a virtual environment using \code{python3 -m venv .venv}. What command must you run in the macOS terminal to activate it? How does the terminal prompt change to show it is active?

    \item \textbf{Dependency Freezing}: What is the exact terminal command to export all currently installed packages and their exact versions into a standard shareable configuration file?

    \item \textbf{Jupyter Kernel State}: You run Cell 1 (\code{x = 5}), then run Cell 2 (\code{x = 10}), then delete Cell 2 completely. What value will Python return if you now execute \code{print(x)} in Cell 3? What should you do to fix this?

    \item \textbf{Microcontroller Verification}: What terminal command checks which USB serial devices (such as an ESP32 or Arduino) are currently recognized by your Mac?
\end{enumerate}

% ==============================================================================
% ANSWERS
% ==============================================================================
\subsection*{Answers}

\begin{enumerate}
    \item \textbf{Environment Isolation}:
    Running \code{sudo pip install} forces packages into the root operating system directory (\code{/usr/bin} and system Python paths). macOS uses this system Python for internal OS automations and security tasks. Overwriting these system libraries with incompatible versions can break macOS desktop tools, Xcode extensions, or system stability.

    \item \textbf{Virtual Environment Mechanics}:
    To activate the environment on macOS, run:
    \begin{equation*}
    \code{source .venv/bin/activate}
    \end{equation*}
    The terminal prompt changes by prepending the environment name in parentheses, such as: \code{(.venv) user@MacBook \~{} \%}.

    \item \textbf{Dependency Freezing}:
    The command is:
    \begin{equation*}
    \code{pip freeze > requirements.txt}
    \end{equation*}

    \item \textbf{Jupyter Kernel State}:
    Python will print \code{10}. Because the background Kernel process maintains state in RAM, deleting a cell does \textbf{not} erase variables already executed in memory. To ensure clean, reproducible execution, you must restart the kernel: select **Kernel $\to$ Restart Kernel and Run All Cells**.

    \item \textbf{Microcontroller Verification}:
    To list connected USB serial devices on macOS, run:
    \begin{equation*}
    \code{ls /dev/cu.usbserial* /dev/cu.wchusbserial*}
    \end{equation*}
    (or simply \code{ls /dev/cu.*} to view all serial interfaces).
\end{enumerate}